\documentclass[11pt]{article}
\usepackage{amsmath,amssymb,amsthm}
\usepackage{hyperref}
\usepackage[margin=1in]{geometry}

\newcommand{\N}{\mathbb{N}}
\newcommand{\Z}{\mathbb{Z}}
\newcommand{\Q}{\mathbb{Q}}
\newcommand{\La}{\Lambda}
\newcommand{\GK}{\operatorname{GK}}
\newcommand{\CPP}{\operatorname{CPP}}
\newcommand{\ALCD}{\operatorname{ALCD}}
\newcommand{\IP}{\operatorname{IP}}
\newcommand{\qpoch}[2]{(#1;#2)_\infty}
\newcommand{\qpochn}[3]{(#1;#2)_{#3}}
\newcommand{\qbin}[2]{\genfrac{[}{]}{0pt}{}{#1}{#2}}
\DeclareMathOperator{\tr}{tr}
\DeclareMathOperator{\wt}{wt}
\DeclareMathOperator{\sgn}{sgn}
\DeclareMathOperator{\Fix}{Fix}
\DeclareMathOperator{\Orb}{Orb}
\DeclareMathOperator{\id}{id}

\theoremstyle{plain}
\newtheorem{theorem}{Theorem}
\newtheorem{conjecture}[theorem]{Conjecture}
\newtheorem{proposition}[theorem]{Proposition}
\newtheorem{lemma}[theorem]{Lemma}
\newtheorem{corollary}[theorem]{Corollary}
\theoremstyle{definition}
\newtheorem{definition}[theorem]{Definition}
\newtheorem{example}[theorem]{Example}
\newtheorem{remark}[theorem]{Remark}

\title{The $Q$-Sch\"utzenberger Identity}
\author{Robin Langer}
\date{\today}

\begin{document}
\maketitle

\begin{abstract}
We define the $Q$-Sch\"utzenberger involution $\psi_Q$ on the
pair set $\mathcal{P}_{a,b,k} = \mathcal{G}_{a+k,c} \times
\mathcal{G}_{b+k,c}$ as the composition of Q-Burge$^{-1}$
(Burge-glue on the first component) with Q-RSK (RSK-peel on
the second component).  This is a sign-preserving,
weight-preserving map that complements the sign-reversing
involution~$\iota$ of~\cite{LOCAL_RULE}.  Together, $\iota$
and~$\psi_Q$ generate a dihedral group action on the pair set,
giving the \emph{$Q$-Sch\"utzenberger identity}: the
fixed-point structure of the combined action constrains the
positivity of the $Q$-kernel $\Phi_{a,b}(t)$.
\end{abstract}

\tableofcontents


%% ============================================================
\section{The classical Sch\"utzenberger involution}
\label{sec:classical}
%% ============================================================

We briefly recall the standard theory to fix notation and
motivate the $Q$-analogue.

\subsection{Two bijections, one involution}

Both the Burge and RSK local rules give bijections between
cylindric plane partitions and pairs (interlacing partition,
ALCD):
\begin{align}
  \psi_B &\colon \CPP(\pi) \;\xrightarrow{\;\sim\;}\;
    \IP \times \ALCD(\pi),
    \label{eq:psi-B} \\
  \psi_R &\colon \CPP(\pi) \;\xrightarrow{\;\sim\;}\;
    \IP \times \ALCD(\pi).
    \label{eq:psi-R}
\end{align}
Both are weight-preserving.  They differ only in the cascade
direction at each face: Burge cascades \textbf{up}, RSK
cascades \textbf{down}.

\begin{definition}[Classical Sch\"utzenberger involution]
\label{def:classical-schutz}
The \emph{Sch\"utzenberger involution} is the composition:
\begin{equation}\label{eq:psi-classical}
  \psi := \psi_R \circ \psi_B^{-1}
  \colon \CPP(\pi) \to \CPP(\pi).
\end{equation}
\end{definition}

\begin{proposition}[Classical properties]
\label{prop:classical}
The map $\psi$ satisfies:
\begin{enumerate}
  \item $\psi^2 = \mathrm{id}$ \textup{(involution)},
  \item $\psi$ preserves total weight
    \textup{(}$|\psi(c)| = |c|$\textup{)},
  \item $\psi$ preserves the base partition~$\gamma$ and the
    ALCD face labels~$\mathfrak{d}$,
  \item the internal vertex partitions change:
    $\psi(c) \neq c$ in general.
\end{enumerate}
\end{proposition}

The involution $\psi$ is a non-trivial symmetry of
$\CPP(\pi)$ that exists purely because there are two valid
local rules for the same algebraic identity
$T_n P_m = \sum_k P_{m-k} T_{n-k}$.


%% ============================================================
\section{The $Q$-Sch\"utzenberger on pairs}
\label{sec:Q-schutz}
%% ============================================================

We now construct the $Q$-analogue on the pair set from
\cite{LOCAL_RULE}, Section~10.

\subsection{Ingredients}

Recall from \cite{LOCAL_RULE}, Sections~12--13:
\begin{itemize}
  \item \textbf{Q-Burge} $\mathfrak{B}$: Burge-peel on the
    first component.  Maps
    $((S, \mathfrak{d}),\, (S', \mathfrak{d}')) \mapsto
    ((S \cup \{\ell_B\}, \mathfrak{d}_B),\,
    (S', \mathfrak{d}'))$.
    Sign-reversing (changes $|S|$ by $+1$).
  \item \textbf{Q-RSK} $\mathfrak{R}$: RSK-peel on the second
    component.  Maps
    $((S, \mathfrak{d}),\, (S', \mathfrak{d}')) \mapsto
    ((S, \mathfrak{d}),\,
    (S' \cup \{\ell_R\}, \mathfrak{d}'_R))$.
    Sign-reversing (changes $|S'|$ by $+1$).
  \item \textbf{Q-Burge$^{-1}$} $= \mathfrak{B}^{-1}$:
    Burge-glue on the first component.  Given $i \in S$,
    removes $i$ from $S$ and inserts into $\mathfrak{d}$ via
    Burge forward (cascade \textbf{up}).
    Sign-reversing (changes $|S|$ by $-1$).
  \item \textbf{Q-RSK$^{-1}$} $= \mathfrak{R}^{-1}$:
    RSK-glue on the second component.  Removes $i' \in S'$
    and inserts into $\mathfrak{d}'$ via RSK forward (cascade
    \textbf{down}).  Sign-reversing.
\end{itemize}

\subsection{The composition}

\begin{definition}[$Q$-Sch\"utzenberger]
\label{def:Q-schutz}
The \emph{$Q$-Sch\"utzenberger map} is the composition:
\begin{equation}\label{eq:psi-Q}
  \boxed{\;
  \psi_Q := \mathfrak{R} \circ \mathfrak{B}^{-1}
  \colon \mathcal{P}_{a,b,k} \to \mathcal{P}_{a,b,k}.
  \;}
\end{equation}
Concretely, given a pair
$p = ((S, \mathfrak{d}),\, (S', \mathfrak{d}'))$ at cascade
level~$k$:
\begin{enumerate}
  \item \textbf{Burge-glue the first component.}\;
    Select the canonical element $i \in S$ (determined by the
    Burge forward rule) and apply Burge-glue:
    \[
      (S, \mathfrak{d})
      \;\xrightarrow{\;\mathfrak{B}^{-1}\;}
      (S \setminus \{i\},\; \tilde{\mathfrak{d}}).
    \]
    The element~$i$ is absorbed into the ALCD via the Burge
    cascade (\textbf{up}).

  \item \textbf{RSK-peel the second component.}\;
    Apply RSK-peel to the (unchanged) second component:
    \[
      (S', \mathfrak{d}')
      \;\xrightarrow{\;\mathfrak{R}\;}
      (S' \cup \{\ell_R\},\; \mathfrak{d}'_R).
    \]
    The element~$\ell_R$ is extracted from the ALCD via the
    RSK cascade (\textbf{up}).
\end{enumerate}
The result is:
\begin{equation}\label{eq:psi-Q-explicit}
  \psi_Q(p) = \bigl(
    (S \setminus \{i\},\, \tilde{\mathfrak{d}}),\;\,
    (S' \cup \{\ell_R\},\, \mathfrak{d}'_R)
  \bigr).
\end{equation}
\end{definition}

\begin{remark}[Both steps cascade \textbf{up}]
\label{rem:both-up}
The Burge-glue uses Burge forward (cascade \textbf{up}) and
the RSK-peel uses RSK inverse (cascade \textbf{up}).  Both
steps push data in the same direction.  This is not a
coincidence: it reflects the fact that
$\mathfrak{B}^{-1}$ inverts the \textbf{down} cascade of
Burge$^{-1}$ by going \textbf{up}, and $\mathfrak{R}$ inverts
the \textbf{down} cascade of RSK forward by going \textbf{up}
via RSK$^{-1}$.
\end{remark}


%% ============================================================
\section{Properties of $\psi_Q$}
\label{sec:properties}
%% ============================================================

\begin{proposition}[Sign and weight preservation]
\label{prop:sign-weight-psi}
The map $\psi_Q$ satisfies:
\begin{enumerate}
  \item \textbf{Sign-preserving:}\;
    $(-1)^{|S(\psi_Q(p))| + |S'(\psi_Q(p))|}
    = (-1)^{|S(p)| + |S'(p)|}$.
  \item \textbf{Weight-preserving:}\;
    $\wt(\psi_Q(p)) = \wt(p)$, where
    $\wt(p) = \sigma(S) + |\mathfrak{d}| + \sigma(S')
    + |\mathfrak{d}'|$.
  \item \textbf{Cascade-preserving:}\;
    $\psi_Q$ maps $\mathcal{P}_{a,b,k}$ to
    $\mathcal{P}_{a,b,k}$ (both components remain at the
    same ground-set level).
\end{enumerate}
\end{proposition}

\begin{proof}
(1) The Burge-glue step changes $|S|$ by $-1$; the RSK-peel
step changes $|S'|$ by $+1$.  Net change in $|S| + |S'|$:
$-1 + 1 = 0$.

(2) Burge-glue preserves the first-component weight
($\sigma(S \setminus \{i\}) + |\tilde{\mathfrak{d}}|
= \sigma(S) + |\mathfrak{d}|$), and RSK-peel preserves the
second-component weight.

(3) Burge-glue maps
$(S, \mathfrak{d}) \in \mathcal{G}_{a+k,c}$ to
$(S \setminus \{i\}, \tilde{\mathfrak{d}})
\in \mathcal{G}_{a+k,c}$ (the ground-set level $a + k$ is
unchanged: the subset shrinks by one element, the ALCD grows
by one level, and $(a+k) - (|S|-1) = (a+k-|S|) + 1$).
Similarly for the second component.
\end{proof}


\subsection{The involution conjecture}

\begin{conjecture}[$Q$-Sch\"utzenberger involution]
\label{conj:psi-involution}
$\psi_Q^2 = \mathrm{id}$ on $\mathcal{P}_{a,b,k}$.
\end{conjecture}

In the classical case, $\psi^2 = \mathrm{id}$ follows from
the Bender--Knuth involution relations: the transposition of
two adjacent rows of the growth diagram, implemented via
Burge and RSK, is self-inverse.  The $Q$-analogue requires
that the composition
\[
  (\mathfrak{R} \circ \mathfrak{B}^{-1})
  \circ
  (\mathfrak{R} \circ \mathfrak{B}^{-1})
  = \mathrm{id},
\]
i.e., that Burge-glue followed by RSK-peel, applied twice,
returns to the original pair.  This amounts to a compatibility
condition between the Burge and RSK cascade directions at the
ground-set level.

\begin{remark}[Sufficient condition]
\label{rem:sufficient}
A sufficient condition for $\psi_Q^2 = \mathrm{id}$ is that
the element~$i$ glued by $\mathfrak{B}^{-1}$ and the
element~$\ell_R$ peeled by $\mathfrak{R}$ satisfy a
``round-trip'' property: applying $\psi_Q$ twice returns
$i$ to~$S$ and removes $\ell_R$ from~$S'$.  In the classical
case, this is guaranteed because the growth-diagram row
transposition is self-inverse.  At the pair level, the
round-trip property needs verification for each profile~$c$.
\end{remark}

\begin{remark}[The $d = 1$ case]
\label{rem:d1}
When $d = 1$, $Q_n = 1$ for all~$n$, the ground set
$\mathcal{G}_{n,c}$ reduces to the standard ALCD set (with
$S = \emptyset$ for the only non-trivial contribution), and
$\psi_Q$ reduces to the classical Sch\"utzenberger involution.
In this case, $\psi_Q^2 = \mathrm{id}$ is the classical
theorem.
\end{remark}


%% ============================================================
\section{The sign problem: why $\psi_Q$ does not prove positivity}
\label{sec:sign-problem}
%% ============================================================

\subsection{The fundamental distinction}

Proving that $Q_{n,c}(q) \in \Z_{\geq 0}[q]$ requires a
\textbf{sign-reversing} involution
$\iota\colon \mathcal{G}_{n,c} \to \mathcal{G}_{n,c}$ that pairs
elements of opposite sign
($|S'| \equiv |S| + 1 \pmod{2}$) while preserving weight.  The
$Q$-Sch\"utzenberger map $\psi_Q$ is \textbf{sign-preserving}
(Proposition~\ref{prop:sign-weight-psi}): it fixes the total
$|S| + |S'|$, and therefore cannot perform the cancellations
that establish positivity.

\begin{center}
\begin{tabular}{l|c|c|c}
\textbf{Map} & \textbf{Sign} & \textbf{Direction} &
  \textbf{Proves positivity?} \\ \hline
$\psi_Q$ ($Q$-Sch\"utzenberger)
  & preserving & horizontal & No \\[3pt]
$\iota$ (peel/glue involution)
  & reversing & vertical & Yes (if canonical peel exists)
\end{tabular}
\end{center}

The directions are:
\begin{itemize}
  \item \textbf{Horizontal} ($\psi_Q$): moves between the
    Burge and RSK realisations at the \emph{same} $|S| + |S'|$.
    This is a change of coordinates---column cascade vs.\ row
    cascade---that does not alter the subsets.
  \item \textbf{Vertical} ($\iota$): moves between subset levels
    $|S| + |S'|$ and $|S| + |S'| \pm 1$, trading subset elements
    for cylindric structure.  This is the sign-changing operation
    that drives cancellation.
\end{itemize}


\subsection{The missing ingredient: the canonical peel}

The sign-reversing involution~$\iota$ requires the
\textbf{canonical peel choice rule}: given
$(S, \mathfrak{d})$, which element~$k$ to peel or glue, and
whether to peel or glue (\cite{LOCAL_RULE}, Conjecture~8.1).
The framework of Sections~\ref{sec:Q-schutz}--\ref{sec:properties}
guarantees that the type signature and weight condition are
satisfied for \emph{any} valid choice of~$k$.  But the choice
must be deterministic, involutive, and must produce fixed points
only at even $|S| + |S'|$.  The $Q$-Sch\"utzenberger
map~$\psi_Q$ does not address this choice.


\subsection{Classical analogy}

The situation mirrors the classical theory.

\begin{center}
\begin{tabular}{l|l}
\textbf{Classical} & \textbf{$Q$-enriched} \\ \hline
Sch\"utzenberger involution on SYT &
  $Q$-Sch\"utzenberger $\psi_Q$ on pairs \\[3pt]
Symmetry (transposes $P$- and $Q$-tableaux) &
  Symmetry (swaps Burge/RSK coordinates) \\[3pt]
Does not prove Littlewood--Richardson rule &
  Does not prove $Q_{n,c} \geq 0$ \\[3pt]
Jeu de taquin / crystal operators prove LR &
  Canonical peel rule proves positivity \\[3pt]
Franklin's sign-reversing involution &
  Sign-reversing $\iota$ for positivity
\end{tabular}
\end{center}

Just as the Sch\"utzenberger involution on standard Young
tableaux is a deep symmetry that does not by itself establish
the Littlewood--Richardson rule, the $Q$-Sch\"utzenberger is a
structural result about the \emph{duality} of the pair-level
framework, not a \emph{cancellation} result.


\subsection{The coordinate-independence conjugacy}

Despite not proving the conjecture, $\psi_Q$ constrains the
search for the canonical peel.

\begin{proposition}[Conjugacy constraint]
\label{prop:conjugacy}
If the canonical peel rule is \emph{coordinate-independent}---depending
only on partition sizes and abstract structure, not on the
column-vs-row resolution of overlaps---then the Burge and RSK
peel/glue involutions are conjugate via~$\psi_Q$:
\begin{equation}\label{eq:conjugacy}
  \mathfrak{V}^R = \psi_Q \circ \mathfrak{V}^B \circ \psi_Q^{-1}.
\end{equation}
\end{proposition}

\begin{proof}
If the canonical choice depends only on $(S, \mathfrak{d})$ at
the abstract level (not on which coordinate system resolves
overlaps), then $\mathfrak{V}^B$ and $\mathfrak{V}^R$ make the
same peel/glue decisions.  They differ only in the face-level
geometry (Burge vs.\ RSK cascade), which is exactly the
transformation performed by~$\psi_Q$.
\end{proof}

\begin{corollary}\label{cor:half-work}
Under the hypothesis of Proposition~\ref{prop:conjugacy}, finding
the canonical peel for Burge automatically gives the RSK version.
The conjugacy~\eqref{eq:conjugacy} also implies
$\psi_Q\colon \Fix(\mathfrak{V}^B) \xrightarrow{\sim}
\Fix(\mathfrak{V}^R)$: the two ``charge decompositions'' are
equivalent.
\end{corollary}


%% ============================================================
\section{The dihedral action}
\label{sec:dihedral}
%% ============================================================

The pair set $\mathcal{P}_{a,b,k}$ now carries two
distinguished maps:
\begin{center}
\begin{tabular}{l|l|l|l}
\textbf{Map} & \textbf{Definition} & \textbf{Sign}
  & \textbf{Order} \\ \hline
$\iota$ & Q-Burge--RSK involution
  (\cite{LOCAL_RULE}, Def.~13.1)
  & reversing & $2$ \\
$\psi_Q$ & $\mathfrak{R} \circ \mathfrak{B}^{-1}$
  (Definition~\ref{def:Q-schutz})
  & preserving & $2$ (conj.)
\end{tabular}
\end{center}

\begin{proposition}[Group structure]
\label{prop:dihedral}
Assuming Conjecture~\ref{conj:psi-involution}, the group
$\Gamma = \langle \iota, \psi_Q \rangle$ generated by the
two involutions is a dihedral group:
\[
  \Gamma \cong D_m
  \quad\text{for some } m \geq 1,
\]
where $m$ is the order of $\iota \circ \psi_Q$.  The group
acts on $\mathcal{P}_{a,b,k}$ preserving weight and with the
sign character:
\[
  \sgn(\gamma) = \begin{cases}
    +1 & \text{if } \gamma \text{ is a product of an even
      number of } \iota\text{'s}, \\
    -1 & \text{if odd}.
  \end{cases}
\]
\end{proposition}

\begin{proof}
The group generated by two involutions $\iota$ and $\psi_Q$
is always dihedral (or cyclic if they commute).  Since $\iota$
is sign-reversing and $\psi_Q$ is sign-preserving, the sign
character distinguishes the two cosets of the rotation subgroup
$\langle \iota \circ \psi_Q \rangle$ in~$\Gamma$.
\end{proof}


\subsection{Orbits and fixed points}

The $\Gamma$-orbits on $\mathcal{P}_{a,b,k}$ decompose into
three types:

\begin{definition}[Orbit types]
\label{def:orbit-types}
Let $\mathcal{O}$ be a $\Gamma$-orbit in
$\mathcal{P}_{a,b,k}$.
\begin{enumerate}
  \item \textbf{Type~I (cancelled).}\;
    $\mathcal{O}$ contains elements of both signs.  These
    contribute zero to $Q_{a+k} Q_{b+k}$ after signed
    summation.
  \item \textbf{Type~II (positive).}\;
    $\mathcal{O}$ consists entirely of elements with sign
    $+1$ (i.e., $|S| + |S'|$ even for all elements).  Every
    element is a fixed point of~$\iota$ (since $\iota$ is
    sign-reversing, it must fix elements that have no
    opposite-sign partner in~$\mathcal{O}$).
  \item \textbf{Type~III (negative).}\;
    $\mathcal{O}$ consists entirely of elements with sign
    $-1$.  Same reasoning: every element is a fixed point
    of~$\iota$.
\end{enumerate}
\end{definition}

\begin{proposition}[Positivity criterion]
\label{prop:positivity-criterion}
$Q_{a+k} \cdot Q_{b+k} \in \Z_{\geq 0}[q]$ if and only if
there are no Type~III orbits at cascade level~$k$.
Equivalently: every fixed point of~$\iota$ in
$\mathcal{P}_{a,b,k}$ has sign $+1$.
\end{proposition}

\begin{proof}
The signed generating function is:
\begin{align*}
  Q_{a+k} \cdot Q_{b+k}
  &= \sum_{p \in \mathcal{P}_{a,b,k}}
    (-1)^{|S(p)| + |S'(p)|}\, q^{\wt(p)} \\
  &= \sum_{\substack{p \in \Fix(\iota)}}
    (-1)^{|S(p)| + |S'(p)|}\, q^{\wt(p)},
\end{align*}
since the non-fixed-point pairs cancel in opposite-sign pairs
under~$\iota$.  This is non-negative if and only if every
fixed point has $|S(p)| + |S'(p)|$ even.
\end{proof}


\subsection{The role of $\psi_Q$}

The $Q$-Sch\"utzenberger $\psi_Q$ acts \emph{within} the
fixed-point set $\Fix(\iota)$ (since $\psi_Q$ preserves sign,
it maps fixed points of $\iota$ to fixed points of $\iota$,
provided $\iota$ and $\psi_Q$ commute on $\Fix(\iota)$).
This gives additional structure:

\begin{proposition}[Refinement of fixed points]
\label{prop:refinement}
If\/ $\psi_Q$ restricts to an involution on $\Fix(\iota)$,
then the fixed-point contribution decomposes as:
\[
  \sum_{p \in \Fix(\iota)} (-1)^{|S|+|S'|}\, q^{\wt(p)}
  \;=\;
  \underbrace{
    \sum_{p \in \Fix(\iota) \cap \Fix(\psi_Q)}
    (-1)^{|S|+|S'|}\, q^{\wt(p)}
  }_{\text{doubly fixed}}
  \;+\;
  \underbrace{
    \sum_{\substack{p \in \Fix(\iota) \\ \psi_Q(p) \neq p}}
    (-1)^{|S|+|S'|}\, q^{\wt(p)}
  }_{\text{$\psi_Q$-paired, same sign}}.
\]
Since $\psi_Q$ is sign-preserving, the $\psi_Q$-paired terms
on the right contribute \emph{positively} (they have the same
sign and pair within the same sign class).  The doubly fixed
points are the irreducible core.
\end{proposition}


%% ============================================================
\section{The $Q$-Sch\"utzenberger identity}
\label{sec:identity}
%% ============================================================

We now state the main identity.

\begin{theorem}[$Q$-Sch\"utzenberger identity]
\label{thm:Q-schutz-identity}
Let $c = (c_0, c_1, c_2)$ be a rank-$3$ profile with
$d = c_0 + c_1 + c_2$, and let $a, b, k \geq 0$.
The $Q$-kernel coefficient $Q_{a+k}(q) \cdot Q_{b+k}(q)$
admits the decomposition:
\begin{equation}\label{eq:Q-schutz-identity}
  \boxed{\;
  Q_{a+k}\, Q_{b+k}
  = \sum_{p \in \Fix(\iota) \cap \Fix(\psi_Q)}
    q^{\wt(p)}
  \;+\;
  2 \sum_{\substack{\{p, \psi_Q(p)\}
    \subseteq \Fix(\iota) \\ \psi_Q(p) \neq p}}
    q^{\wt(p)},
  \;}
\end{equation}
where the second sum runs over unordered $\psi_Q$-orbits of
size~$2$ within $\Fix(\iota)$, and both sums contribute
non-negative coefficients (all terms have
$|S(p)| + |S'(p)|$ even, hence sign~$+1$).
\end{theorem}

\begin{proof}
By Proposition~\ref{prop:positivity-criterion}, the signed sum
reduces to $\Fix(\iota)$.  By
Proposition~\ref{prop:refinement}, this decomposes into
doubly-fixed points and $\psi_Q$-pairs within $\Fix(\iota)$.
Since $\psi_Q$ is weight-preserving and sign-preserving, each
$\psi_Q$-pair $\{p, \psi_Q(p)\}$ contributes
$2\, q^{\wt(p)}$ (both elements have the same weight and sign
$+1$).  The doubly-fixed points contribute
$1 \cdot q^{\wt(p)}$ each.
\end{proof}

\begin{remark}[Conditional on two conjectures]
\label{rem:conditional}
Theorem~\ref{thm:Q-schutz-identity} is conditional on:
\begin{enumerate}
  \item[(a)] The existence of the sign-reversing involution
    $\iota$ with only positive fixed points (Warnaar's
    Conjecture~2.7 at the pair level), and
  \item[(b)] Conjecture~\ref{conj:psi-involution}
    ($\psi_Q^2 = \mathrm{id}$).
\end{enumerate}
Point~(a) is the hard conjecture.  Point~(b) is a structural
property of the Burge--RSK interaction at the ground-set level,
which should follow from the classical Bender--Knuth theory
adapted to the $Q$-enriched setting.
\end{remark}


\subsection{The identity at $b = 0$}

Setting $b = 0$ in~\eqref{eq:Q-schutz-identity} gives a
formula for $Q_a$ itself, since $Q_0 = 1$:
\begin{equation}\label{eq:b0}
  Q_{a+k}
  = Q_{a+k} \cdot Q_k \big|_{b=0}
  = Q_{a+k} \cdot 1.
\end{equation}
More precisely, the pair set at $b = 0$ is
$\mathcal{G}_{a+k} \times \mathcal{G}_{k}$, and
$\mathcal{G}_0 = \{(\emptyset, \mathfrak{d}_0)\}$ is a
singleton.  The second component is trivially fixed by both
$\iota$ and $\psi_Q$, so:
\begin{equation}\label{eq:Q-from-Phi}
  Q_a = [t^0]\, \Phi_{a,0}(t)
  = \sum_{\substack{p \in \Fix(\iota) \cap \Fix(\psi_Q)
    \\ p \in \mathcal{G}_a \times \mathcal{G}_0}}
    q^{\wt(p)}
  \;+\;
  2 \sum_{\substack{\{p, \psi_Q(p)\}
    \subseteq \Fix(\iota) \cap (\mathcal{G}_a \times
    \mathcal{G}_0)}}
    q^{\wt(p)}.
\end{equation}
Since the second component is the singleton
$(\emptyset, \mathfrak{d}_0)$, the pair reduces to a single
ground-set element $(S, \mathfrak{d}) \in \mathcal{G}_{a,c}$,
and the $Q$-Sch\"utzenberger identity becomes a statement
about the ground set alone:
\begin{equation}\label{eq:Q-ground-set-identity}
  Q_a = \sum_{(S, \mathfrak{d}) \in \Fix(\iota_a)}
    (-1)^{|S|}\, q^{\sigma(S) + |\mathfrak{d}|}
  \;\in\; \Z_{\geq 0}[q],
\end{equation}
where $\iota_a$ is the restriction of~$\iota$ to
$\mathcal{G}_{a,c} \times \{(\emptyset, \mathfrak{d}_0)\}
\cong \mathcal{G}_{a,c}$.


%% ============================================================
\section{Special cases}
\label{sec:special}
%% ============================================================

\subsection{The case $d = 1$}

When $d = 1$ (one box per period), $Q_n = 1$ for all $n$, the
ground set $\mathcal{G}_{n,c}$ has a single element
$(\emptyset, \mathfrak{d})$ for each ALCD~$\mathfrak{d}$ (the
Euler-factor subsets $S \neq \emptyset$ contribute
cancelling pairs that sum to zero).  The $Q$-Sch\"utzenberger
$\psi_Q$ reduces to the classical Sch\"utzenberger involution,
and the identity~\eqref{eq:Q-schutz-identity} becomes:
\[
  Q_{a+k} \cdot Q_{b+k} = 1
  = |\ALCD_{a+k}(c) \times \ALCD_{b+k}(c)|\big|_{\text{signed}}.
\]

\subsection{The case $d = 2$ (Rogers--Ramanujan)}

For $d = 2$ with profile $c = (2, 0, 0)$:
$Q_n = q^{n^2}$ (first Rogers--Ramanujan identity).  The
product is:
\[
  Q_{a+k} \cdot Q_{b+k} = q^{(a+k)^2 + (b+k)^2}.
\]
This is a single monomial, so the pair set at each cascade
level has exactly one doubly-fixed point contributing
$q^{(a+k)^2 + (b+k)^2}$, and all other elements cancel
under~$\iota$.  The $Q$-Sch\"utzenberger acts trivially on the
unique fixed point ($\psi_Q = \mathrm{id}$ on $\Fix(\iota)$).

For $c = (1, 1, 0)$: $Q_n = q^{n(n+1)}$ (second
Rogers--Ramanujan), and the same monomial structure holds:
\[
  Q_{a+k} \cdot Q_{b+k} = q^{(a+k)(a+k+1) + (b+k)(b+k+1)}.
\]


\subsection{The case $d = 4$, $c = (2,2,0)$}

Here $Q_0 = 1$, $Q_1 = q^2 + q^3 + q^4 + q^5$.  The product
$Q_1 \cdot Q_1 = (q^2 + q^3 + q^4 + q^5)^2$ has multiple
terms, so the doubly-fixed points and $\psi_Q$-pairs of
$\Fix(\iota)$ contribute a non-trivial decomposition.  The
$\psi_Q$-orbits of size~$2$ within $\Fix(\iota)$ account for
the ``extra factor of~$2$'' in the coefficients of
$Q_1^2 = q^4 + 2q^5 + 3q^6 + 4q^7 + 3q^8 + 2q^9 + q^{10}$.


%% ============================================================
\section{Relation to the shifted factorisation}
\label{sec:shifted}
%% ============================================================

The shifted factorisation (\cite{ADJOINT}, Part~III):
\begin{equation}\label{eq:shifted}
  \Phi_{a,b}(t) = Q_a\, Q_b + t\, \Phi_{a+1,b+1}(t)
\end{equation}
acquires a refined interpretation through the dihedral action.

At cascade level $k = 0$, the pair set is
$\mathcal{P}_{a,b,0} = \mathcal{G}_a \times \mathcal{G}_b$.
The $\Gamma$-action decomposes it into orbits, and the
fixed-point contribution is $Q_a \cdot Q_b$.

At level $k = 1$, the pair set is
$\mathcal{G}_{a+1} \times \mathcal{G}_{b+1}$, and the
fixed-point contribution is $Q_{a+1} \cdot Q_{b+1}$.

The shifted factorisation says: the $k = 0$ contribution and
the $k \geq 1$ contributions assemble into $\Phi_{a,b}(t)$,
with each level independently non-negative (assuming the
conjecture).  The $Q$-Sch\"utzenberger provides an
\emph{intra-level symmetry}: within each cascade level~$k$,
the fixed points of $\iota$ are further organised by
$\psi_Q$ into doubly-fixed points and $\psi_Q$-pairs.

\begin{remark}[Inter-level vs.\ intra-level]
\label{rem:inter-intra}
The cascade-coupled involution~$\iota$ can pair elements at
\emph{different} cascade levels (Definition~11.2 of
\cite{LOCAL_RULE}).  The $Q$-Sch\"utzenberger~$\psi_Q$
operates \emph{within} a single level.  Together they provide
both inter-level cancellation (via $\iota$) and intra-level
symmetry (via $\psi_Q$).
\end{remark}


%% ============================================================
\section{The $Q$-Sch\"utzenberger identity: summary}
\label{sec:summary}
%% ============================================================

We summarise the logical structure.

\begin{center}
\begin{tabular}{r|l|l}
\textbf{Map} & \textbf{Type} & \textbf{Role} \\ \hline
$\iota$ & sign-reversing involution &
  cancels opposite-sign pairs \\
$\psi_Q$ & sign-preserving involution (conj.) &
  organises fixed points of $\iota$ \\
$\Gamma = \langle \iota, \psi_Q \rangle$ &
  dihedral group & full symmetry of pair set
\end{tabular}
\end{center}

\medskip
\noindent\textbf{The identity.}\;
The $Q$-kernel coefficient decomposes as:
\[
  Q_{a+k}\, Q_{b+k}
  \;=\;
  \underbrace{|\Fix(\iota) \cap \Fix(\psi_Q)|_q}_{%
    \text{doubly fixed}}
  \;+\;
  \underbrace{2\, |\text{$\psi_Q$-pairs in } \Fix(\iota)|_q}_{%
    \text{Sch\"utzenberger pairs}},
\]
where $|\cdot|_q$ denotes the $q$-weighted count.  All terms
are non-negative.

\medskip
\noindent\textbf{The programme.}
\begin{enumerate}
  \item Construct the sign-reversing involution~$\iota$
    (the hard step --- Warnaar's conjecture).
  \item Verify $\psi_Q^2 = \mathrm{id}$ (the
    $Q$-Bender--Knuth relations).
  \item Characterise the doubly-fixed points
    $\Fix(\iota) \cap \Fix(\psi_Q)$ --- these are the
    irreducible positivity core.
  \item At $b = 0$: the identity reduces to a formula for
    $Q_a$ in terms of the ground set $\mathcal{G}_{a,c}$
    alone (equation~\eqref{eq:Q-ground-set-identity}).
\end{enumerate}


%% ============================================================
\section{The orthogonal structure of the problem}
\label{sec:orthogonal}
%% ============================================================

The analysis reveals that the positivity problem decomposes into
two \emph{orthogonal} components.

\begin{center}
\begin{tabular}{l|l|l}
& \textbf{Horizontal (symmetry)} &
  \textbf{Vertical (cancellation)} \\ \hline
\textbf{Map} &
  $\psi_Q$ ($Q$-Sch\"utzenberger) &
  $\iota$ (peel/glue) \\[3pt]
\textbf{Action} &
  Swaps Burge $\leftrightarrow$ RSK &
  Trades $S \leftrightarrow \mathfrak{d}$ \\[3pt]
\textbf{Subsets $S, S'$} &
  Fixed &
  Change by $\pm 1$ \\[3pt]
\textbf{ALCDs $\mathfrak{d}, \mathfrak{d}'$} &
  Change (coord.\ system) &
  Change (peel/glue) \\[3pt]
\textbf{Sign} &
  Preserved &
  Reversed \\[3pt]
\textbf{Status} &
  \textbf{Defined} (this note) &
  \textbf{Conjectural} (\cite{LOCAL_RULE}) \\[3pt]
\textbf{Role} &
  Constrains the problem &
  Solves the problem
\end{tabular}
\end{center}

The horizontal component ($\psi_Q$) tells us the \emph{stage is
well-built}: the pair-level framework supports a clean duality
between two coordinate systems, confirming that this is the
correct level of abstraction for the $Q$-enriched theory.

The vertical component ($\iota$) is the \emph{actor}: it must
perform the sign-reversing cancellation that establishes
positivity.  Its construction requires the canonical peel rule,
which remains open.

The two components interact through the dihedral group
$\Gamma = \langle \iota, \psi_Q \rangle$
(Section~\ref{sec:dihedral}), and through the conjugacy
$\mathfrak{V}^R = \psi_Q \circ \mathfrak{V}^B \circ \psi_Q^{-1}$
(Proposition~\ref{prop:conjugacy}).  But they serve fundamentally
different purposes.


%% ============================================================
\section{Status of the programme}
\label{sec:status}
%% ============================================================

\subsection{What is established}

\begin{enumerate}
  \item \textbf{Ground set.}\;
    $\mathcal{G}_{n,c}$ with signed generating function
    $Q_{n,c}(q)$ (\cite{UNDERLYING_SET, FIRST_ATTEMPT}).

  \item \textbf{Two local rules.}\;
    Burge and RSK peel/glue on pairs, with shared type
    signature, same face label, same weight condition
    (\cite{Q_LOCAL}).

  \item \textbf{$Q$-Sch\"utzenberger.}\;
    $\psi_Q = \mathfrak{R} \circ \mathfrak{B}^{-1}$
    well-defined, weight-preserving, sign-preserving
    (this note).

  \item \textbf{Locality principle.}\;
    Effective operations in the $Q$-enriched setting must be
    local---one face at a time---matching the shifted
    factorisation $\Phi_{a,b}(t) = Q_a Q_b + t\,\Phi_{a+1,b+1}(t)$
    (\cite{LOCAL_RULE}, Theorem~7.1).

  \item \textbf{Propagation obstruction.}\;
    The global Sch\"utzenberger involution fails in the
    $Q$-enriched setting (\cite{LOCAL_RULE}, Section~4),
    confirming that global approaches do not work.

  \item \textbf{Cyclic descent injection.}\;
    For $d \not\equiv 0 \pmod{3}$, the max-$1$ case has a valid
    injection establishing monotonicity of $h_{1,c}(q)$
    coefficients (\cite{NOTES}, Section~8).

  \item \textbf{Computational verification.}\;
    $Q_{n,c}(q) \in \Z_{\geq 0}[q]$ verified for
    $d = 1, 2, 4, 5, 7$ (all tested profiles pass).
\end{enumerate}


\subsection{What remains open}

\begin{enumerate}
  \item \textbf{Canonical peel rule}
    (\cite{LOCAL_RULE}, Conjecture~8.1).\;
    The deterministic choice of which element~$k$ to peel or
    glue, given $(S, \mathfrak{d})$.  This is the central open
    problem.

  \item \textbf{$Q$-Bender--Knuth relations.}\;
    Conjecture~\ref{conj:psi-involution}: $\psi_Q^2 = \id$.
    Structural property of Burge--RSK interaction at the
    ground-set level.

  \item \textbf{Fixed-point characterisation.}\;
    Identify the ``irreducible'' pairs that are fixed by both
    $\iota$ and $\psi_Q$.  Their cardinality at $b = 0$ must be
    $Q_{n,c}(1) = \bigl(\frac{(d+1)(d+2)}{6} - 1\bigr)^n$.

  \item \textbf{Charge statistic.}\;
    The weight function on fixed points defines the ``missing
    charge.''  Understanding its structure (analogue of
    Lascoux--Sch\"utzenberger charge?) is a further goal.

  \item \textbf{The $d \equiv 0 \pmod{3}$ case.}\;
    The cyclic descent injection fails here (balanced slack has
    no cyclic descent).  A different mechanism is needed, possibly
    involving the pair-based framework at deeper ALCD levels.
\end{enumerate}


\subsection{Where to look for the canonical peel}

The $Q$-Sch\"utzenberger, while not proving the conjecture,
narrows the search in three ways.

\paragraph{1.\ Coordinate independence.}
By Proposition~\ref{prop:conjugacy}, a coordinate-independent
peel rule for Burge automatically gives the RSK version.  This
suggests looking for a rule based on \emph{abstract} properties:
sizes of horizontal strips, face-label magnitudes, or diagonal
weights---rather than specific column or row indices.

\paragraph{2.\ The shifted factorisation as guide.}
The recurrence
$\Phi_{a,b}(t) = Q_a Q_b + t\,\Phi_{a+1,b+1}(t)$ peels one
layer at a time (\cite{LOCAL_RULE}, Section~7.2).  The canonical
peel should correspond to the $k = 1 \to k = 0$ reduction:
peel the face whose removal corresponds to stripping the leading
term from the $Q$-kernel.

\paragraph{3.\ The cyclic descent as prototype.}
For the max-$1$ case ($d \not\equiv 0 \pmod{3}$), the cyclic
descent injection (\cite{NOTES}, Section~8) provides an explicit
peel rule based on the slack vector.  This rule is
coordinate-independent and produces an injection from
weight-$(w-1)$ to weight-$w$ configurations.  Extending this to
general ALCD depth is the key challenge.


%% ============================================================
\section{Open questions}
\label{sec:open}
%% ============================================================

\begin{enumerate}
  \item \textbf{$Q$-Bender--Knuth.}\;
    Does $\psi_Q^2 = \mathrm{id}$?  Is there a
    $Q$-analogue of the Bender--Knuth involution that
    mediates between the Burge and RSK cascades at the
    ground-set level?

  \item \textbf{Order of $\iota \circ \psi_Q$.}\;
    What is the order~$m$ of the dihedral group
    $\Gamma = \langle \iota, \psi_Q \rangle$?  In the
    classical case, $\psi$ has order~$2$ and
    $\iota \circ \psi$ may have larger order.  Does the
    $Q$-enrichment change the group structure?

  \item \textbf{Doubly-fixed-point formula.}\;
    Is there a closed form for the $q$-count of
    $\Fix(\iota) \cap \Fix(\psi_Q)$?  In the $d = 2$ case,
    this is a single monomial $q^{(a+k)^2 + (b+k)^2}$
    (or $q^{(a+k)(a+k+1) + (b+k)(b+k+1)}$).  Is there a
    quadratic-form structure for general~$d$?

  \item \textbf{Crystal-base interpretation.}\;
    The classical Sch\"utzenberger involution has an
    interpretation via crystal bases for
    $\widehat{\mathfrak{sl}}_r$
    (cf.~\cite{COMMUTATION}).  Does the
    $Q$-Sch\"utzenberger correspond to a crystal operator
    on the $Q$-enriched pair set?

  \item \textbf{Computational verification.}\;
    For small $d$ and $n$, enumerate $\mathcal{P}_{a,b,k}$,
    construct $\iota$ and $\psi_Q$ explicitly, verify the
    identity~\eqref{eq:Q-schutz-identity}, and tabulate the
    doubly-fixed-point counts.
\end{enumerate}


\begin{thebibliography}{99}
\bibitem{ADJOINT}
  R.~Langer,
  \emph{The $Q$-Operator and its Adjoint: From Level~3 Operator
  Traces to Cascade Commutators},
  companion note, 2026.

\bibitem{Borodin}
  A.~Borodin,
  \emph{Periodic Schur process and cylindric partitions},
  Duke Math.~J.\ \textbf{140} (2007), 391--468.

\bibitem{COMMUTATION}
  R.~Langer,
  \emph{From Commutation Relations to Growth Diagrams:
  A Summary of the Operator--Bijection Correspondence},
  companion note.

\bibitem{FIRST_ATTEMPT}
  R.~Langer,
  \emph{Towards an Involution Proof of Positivity for
  $Q_{n,c}(q)$},
  companion note.

\bibitem{LOCAL_RULE}
  R.~Langer,
  \emph{The Generalized Local Rule and the Burge--RSK
  Involution},
  companion note, 2026.

\bibitem{Q_LOCAL}
  R.~Langer,
  \emph{Two Equivalent Local Rules on Ground-Set Pairs},
  companion note, 2026.

\bibitem{NOTES}
  R.~Langer,
  \emph{Weighted Local Rules and the Cyclic Descent
  Injection},
  companion note.

\bibitem{UNDERLYING_SET}
  R.~Langer,
  \emph{The Underlying Set: Cylindric Plane Partitions, ALCDs,
  and the Ground Set for the Involution},
  companion note.

\bibitem{langer-thesis}
  R.~Langer,
  \emph{Cylindric Partitions and Bijective Proofs},
  PhD thesis, University of Melbourne.

\bibitem{vanLeeuwen}
  M.~A.~A.~van~Leeuwen,
  \emph{The Robinson--Schensted and Sch\"utzenberger
  algorithms, an elementary approach},
  Electron.~J.~Combin.\ \textbf{3} (1996), R15.
\end{thebibliography}

\end{document}
